\section{Install VS Code, the program you will write in}

VS Code is a free text editor from Microsoft. Go to \texttt{code.visualstudio.com}, press the large download button, open the file you receive, and accept the choices the installer offers. On Windows, leave the box marked \textbf{Add to PATH} ticked, because that lets the terminal find it.

\begin{figure}[H]\centering\includegraphics[width=\linewidth]{web_vscode.png}\caption{The VS Code download page. The site changes its wording from time to time, but the download button stays near the top. If a cookie notice appears, either choice is fine.}\end{figure}

When VS Code first opens it may offer to sign you in to an assistant or to choose a theme. You can close these offers, since nothing in this guide needs them.

\section{Install Git, the program that keeps your history}

Git records each saved version of your files, so you can always go back. Go to \texttt{git-scm.com/downloads} and choose your operating system.

\begin{figure}[H]\centering\includegraphics[width=\linewidth]{web_git.png}\caption{The Git download page. Choose Windows, macOS or Linux, then follow the instructions on the page that opens.}\end{figure}

On Windows, run the installer and press \textbf{Next} on every screen, as the suggested settings are fine. On a Mac, open the Terminal application and type \texttt{git --version}, and the system will offer to install it. On Linux, use your system's package manager, for example \texttt{sudo apt install git}.

Close VS Code and open it again so that it notices Git. Then open a terminal and tell Git who you are. Git writes this name and email on every version you save, and they become public if you publish your work.

\begin{lstlisting}[language=bash]
git config --global user.name "Your Name"
git config --global user.email "you@example.org"
git config --global init.defaultBranch main
\end{lstlisting}

These commands print nothing when they succeed. To check, type \texttt{git --version}, which should print a version number.

\section{Install Python}

Python is a programming language, and here it runs small scripts. Go to \texttt{python.org/downloads} and press the yellow download button.

\begin{figure}[H]\centering\includegraphics[width=\linewidth]{web_python.png}\caption{The Python download page. The page recognises your operating system and offers the right version.}\end{figure}

On Windows, tick \textbf{Add python.exe to PATH} on the first screen of the installer before you press \textbf{Install Now}. It is easy to miss, and without it the terminal cannot find Python. Close and reopen VS Code, then check that it worked:

\begin{lstlisting}[language=bash]
python --version
\end{lstlisting}

This should print a version number such as \texttt{Python 3.13.0}. On a Mac or on Linux, the command may be \texttt{python3} instead of \texttt{python}. Use whichever works, wherever this guide says \texttt{python}.

\section{Create a GitHub account}

GitHub is a website that stores copies of Git projects and lets other people see them. Git is the program on your computer, and GitHub is the place online.

Go to \texttt{github.com/signup}. Enter your email, choose a password and pick a username. Choose the username with care, because it appears in your web addresses and is awkward to change later. Your real name, or a close form of it, works well for professional use. GitHub will send a code to your email to confirm the account. The free plan is all you need, so if you are offered a paid plan or asked to fill in a survey, you can skip it.

%SHOT{in_signup.png}{The GitHub sign-up page.}

The first time you send work to GitHub, a window will ask you to sign in through your browser. Accept it, and your computer will remember the permission.
